\chapter{The Category of GSLTs}
\label{ch:catgslt}

Part~III consists of constructions that take a theory and produce a theory. To say
anything about such constructions --- that they are functorial, that they are monads,
that they have adjoints --- we need a category of theories. This section builds it, and
the build is less routine than one might expect: the obvious definition is wrong, and
the way it is wrong is instructive.

\section{The naive definition and three defects}

The obvious move is to take objects to be GSLTs and morphisms to be signature-preserving
translations that preserve bisimulation. Three things go wrong.

\paragraph{The constant map.} Send every term to $0$. This preserves signatures
trivially and preserves bisimulation trivially, since the image has no transitions to
fail to match. It is a morphism from anything to anything, and its presence means the
category orders nothing.

\paragraph{Milner's encoding is not a map of terms.} The canonical encoding of the
$\lambda$-calculus into the $\pi$-calculus~\cite{Milner1992} --- the very construction
that established the modern practice of comparing calculi --- sends a $\lambda$-term not
to a $\pi$-term but to a $\pi$-term \emph{parameterised by a name}: $\sem{M}_u$, the
process that realises $M$ and reports its result on $u$. Signature preservation does not
typecheck, because $\mathrm{App}$ is not sent to any single $\pi$ operator; it is sent to
a construction involving restriction, parallel composition, and output. Requiring
morphisms to be signature homomorphisms excludes the motivating example.

\paragraph{Ambient contexts observe too much.} Suppose one encodes $\pi$ into rho. A rho
context may quote a process, compare quotations, and generally observe the syntactic
form of what it is handed --- capabilities no $\pi$ context has. Bisimulation computed
over \emph{all} rho contexts will therefore separate the images of $\pi$-terms that no
$\pi$ context can separate. Under the naive definition, no interesting encoding is ever
bisimulation-preserving, for a reason that has nothing to do with the encoding.

\section{Contexts as the primary structure}

The three defects have one cure. The first says we need a non-degeneracy condition; the
second says terms are the wrong primitive; the third says the observer must be
restricted to what the source can express. Taking contexts rather than terms as primary
addresses all three at once, because a term is a nullary context, and because the
observers of a term are exactly the contexts one can place around it.

A theory presents a symmetric multicategory of contexts. Its objects are
\emph{interfaces}, recording the sort of a hole, its binding stage, and its interaction
surface. Its multimorphisms are contexts with holes of the given interfaces and a result
of the given interface; composition is plugging. Terms are the nullary contexts.

Transitions are then labelled by contexts: rather than $P \to P'$, we record
$P \xrightarrow{\;C\;} P'$, meaning that placing $P$ in the context $C$ enables the step.
Bisimulation is taken over these context-labelled transitions.

\begin{definition}[Morphism of theories]
\label{ob:def:morphism}
A morphism of GSLTs is a pseudofunctor of context multicategories that is
bisimulation-preserving for context-labelled transitions, where the bisimulation on the
target is computed only over the image of the source's contexts.
\end{definition}

Each clause repairs one defect. Because the map is defined on contexts and terms are
nullary contexts, an encoding like Milner's --- which sends a term to a
context-parameterised process --- is expressible: the parameterisation is the interface.
Because bisimulation on the target is computed over the image, the ambient-observation
defect disappears by construction; the encoded $\pi$-terms are compared by the encoded
$\pi$ contexts, which is the comparison one meant to make. And the constant map is
excluded by the non-degeneracy conditions of the next subsection.

\begin{definition}[The category $\catGSLT$]
\label{ob:def:GSLTcat}
$\catGSLT$ is the category whose $0$-cells are GSLTs and whose morphisms are the maps of
Definition~\ref{ob:def:morphism}: bisimulation-preserving for the context-labelled
transition relation, which is thereby a congruence.
\end{definition}

This is the official definition for the remainder of the document. The additional
conditions imposed by particular constructions --- quote-faithfulness for the cost
functor of Chapter~\ref{ch:costmonad}, preservation of $\Pr$ and $\rsq$ for the type-system
functor of Chapter~\ref{ch:hypercube} --- are stated where those constructions need
them, and are refinements of Definition~\ref{ob:def:morphism} rather than rivals to it.

\section{Hosting and exhausting}

Two conditions prevent the comparison from being trivial.

\begin{definition}[Hosting]
A morphism $F : G \to H$ is \emph{hosting} (faithful) when the target can run the
source: the encoding reflects as well as preserves the context-labelled transition
structure, so that no branching present in $G$ is lost in $H$.
\end{definition}

\begin{definition}[Exhausting]
$F : G \to H$ is \emph{exhausting} (dense) when the target is nothing the source cannot
assemble: every context of $H$ relevant to the image is, up to the equivalence, in the
image of a context of $G$.
\end{definition}

The constant map fails hosting immediately. Together the conditions induce a preorder on
theories whose degrees \emph{refine} Turing completeness rather than replacing it: the
classical notion is recovered as the trace collapse, obtained by coarsening
context-labelled bisimulation all the way to the input/output relation.

\begin{remark}[Why not Turing completeness]
Expressiveness comparisons between computational models are conventionally anchored to
Turing completeness. For functions this works. For interactive systems it does not,
because Turing completeness is a statement about the input/output relation, and the
input/output relation is the coarsest invariant in the linear-time/branching-time
spectrum~\cite{vanGlabbeek1990}, while the equivalences that matter for interaction sit
at the other end. Calculi of identical Turing power are routinely separated by
uniform-encoding criteria, so the yardstick is measuring something other than what is
being compared.

The obvious repair --- demand behavioural rather than computational universality ---
merely relocates the problem, since the arbiter of which behavioural equivalence to use
is an adequate modal logic, and the logical apparatus available in the ambient setting is
inherited from the $\lambda$-calculus. An absolute notion is quietly $\lambda$-relative.
Definition~\ref{ob:def:morphism} replaces the yardstick with a relative property: fix a
probe and ask what that probe can see.
\end{remark}

\begin{remark}[A note on complexity]
The same preference for intensional over extensional comparison has a consequence for
complexity theory that we record without pursuing. Equivalence of regular languages is
\textsc{PSPACE}-complete, while the corresponding bisimulation problem is polynomial;
equivalence of context-free languages is undecidable, while the corresponding
bisimulation problem is tractable. We want the witnesses, not the extensions they give
rise to. Infinitary objects like languages cannot be grasped directly, so an intensional
representation is needed --- and once one has an intensional representation, one has a
bisimulation problem rather than a language-comparison problem.
\end{remark}


The three sections that follow have the same shape. Each takes a theory presented as an
interaction cut and produces another such theory, decorated with apparatus. Each is a
monad. Each resolves through a free/forgetful adjunction that installs and removes the
apparatus. And in each case, when the base theory is sufficiently expressive, the
apparatus can alternatively be encoded back into the base --- so that the base performs
the metering, or the logging, or the checking, with its own computation. We note that
possibility where it arises and do not develop it.
